% TOP_PROBLEM: {"id": 1194, "title": "Existence of Generalized Moonshine", "category_id": 4, "category": {"id": 4, "name": "algebra", "display_name": "Algebra", "description": "Group theory, ring theory, field theory, and algebraic structures.", "slug": "algebra", "order_index": 4, "created_at": "2026-07-31T15:26:25.671Z"}, "status": "open"}
% FIRST_POSTED: null
% ENTRY_KIND: statement_only
% Imported from: batch04.tex; UnsolvedMath ID 1194
\documentclass[11pt]{article}
\usepackage[margin=25mm]{geometry}
\usepackage{fontspec}
\setmainfont{Latin Modern Roman}
\usepackage{amsmath,amssymb,mathtools}
\usepackage{unicode-math}
\setmathfont{Latin Modern Math}
\usepackage{xurl,hyperref}
\hypersetup{colorlinks=true,urlcolor=blue}
\setcounter{secnumdepth}{0}
\setlength{\parindent}{0pt}
\setlength{\parskip}{5pt}
\setlength{\emergencystretch}{3em}
\newcommand{\sref}[2]{\hyperlink{source:#1:#2}{[S#2]}}
\newcommand{\cataloguescope}[1]{\paragraph{Recorded target.}#1}
\begin{document}
% UnsolvedMath ID 1194; source catalogue code ALG-030
\setcounter{section}{1193}
\section{Existence of Generalized Moonshine}\label{1194}
{\small\sffamily UnsolvedMath ID 1194 · Algebra}
\subsection{Definitions and mathematical statement}

\paragraph{Shared notation.}$\mathbb N=\{1,2,\ldots\}$, and $\mathbb Z,\mathbb Q,\mathbb R,\mathbb C$ denote the integers, rationals, reals and complex numbers. Any other conventions are specified locally.


Let \(\mathbb M\) be the Monster finite simple group, and let
\(C_{\mathbb M}(g)=\{h\in\mathbb M:hg=gh\}\).
Put \(\mathfrak H=\{\tau\in\mathbb C:\operatorname{Im}\tau>0\}\) and
\(q=e^{2\pi i\tau}\).
A graded projective representation of \(C_{\mathbb M}(g)\) is a
\(\mathbb Q\)-graded complex vector space
\(V(g)=\bigoplus_{r\in\mathbb Q}V(g)_r\), with finite-dimensional graded
pieces, on which the centralizer acts by grading-preserving linear
maps specified up to nonzero scalar.
A lift \(\widetilde h\) chooses a representative linear map for \(h\).
A genus-zero congruence group means a discrete subgroup
\(\Gamma\le\mathrm{SL}_2(\mathbb R)\) containing a principal congruence
subgroup of \(\mathrm{SL}_2(\mathbb Z)\), whose compactified modular
curve has genus zero.  A Hauptmodul is a meromorphic function
generating the function field of that modular curve.
Let \(J(\tau)=j(\tau)-744=q^{-1}+196884q+\cdots\) be the normalized
classical modular invariant.
Two holomorphic functions are proportional below if one is a
nonzero constant multiple of the other.
\paragraph{Statement recorded in the supplied source.}
Do there exist, for every \(g\in\mathbb M\), a graded projective
representation \(V(g)\) of \(C_{\mathbb M}(g)\), and, for every commuting
pair \(g,h\in\mathbb M\), a holomorphic function \(Z(g,h,\tau)\) on
\(\mathfrak H\), satisfying all of Norton's conditions?
\begin{enumerate}
\item For a suitable lift \(\widetilde h\),
\[
 Z(g,h,\tau)=\sum_{r\in\mathbb Q}
       \operatorname{Tr}(\widetilde h|V(g)_r)\,q^{r-1},
\]
with convergence to the indicated holomorphic function.
\item For every \(x\in\mathbb M\), the functions
\(Z(xgx^{-1},xhx^{-1},\tau)\) and \(Z(g,h,\tau)\) are proportional.
\item Each \(Z(g,h,\tau)\) is either constant or a Hauptmodul for a
genus-zero congruence group.
\item For every
\(\begin{pmatrix}a&b\\c&d\end{pmatrix}\in\mathrm{SL}_2(\mathbb Z)\),
\[
 Z(g^ah^c,g^bh^d,\tau)
   \quad\text{is proportional to}\quad
 Z\!\left(g,h,\frac{a\tau+b}{c\tau+d}\right).
\]
\item \(Z(g,h,\tau)=J(\tau)\) if and only if \(g=h=1\).
\end{enumerate}
These conditions state generalized moonshine for all commuting pairs,
not merely the untwisted character associated with one Monster element.
The cited work proves the historical conjecture. \sref{1194}{2}
\subsection{Short English statement}
Do graded projective centralizer representations and compatible trace functions exist for every commuting pair in the Monster, with conjugation and modular covariance, the constant-or-Hauptmodul property, and the prescribed identity normalization? Carnahan's cited work establishes this historical conjecture.
\subsection{Sources}
\begin{description}
\item[\hypertarget{source:1194:1}{S1}] ULAM.AI and the UnsolvedMath Contributors, \emph{UnsolvedMath: A Curated Collection of Open Mathematics Problems}, record 1194 (ALG-030). CC BY 4.0. \url{https://huggingface.co/datasets/ulamai/UnsolvedMath}.
\item[\hypertarget{source:1194:2}{S2}] Scott Huai-Lei Carnahan, \emph{Generalized Moonshine IV: Monstrous Lie algebras}, arXiv:1208.6254, revised 2016, Section 1.1, Norton's conjecture with conditions (1)--(5). \url{https://arxiv.org/abs/1208.6254}.
\end{description}
\label{end:1194}
\end{document}
